\section{Cycle Properties}

In this section, we prove and verify the main properties of Cycles. Each
property is stated mathematically, then shown to hold via a corresponding
verified lemma in Scala using the Stainless system.

\begin{itemize}
  \item Element access: \texttt{cycle(key) == cycle.values(mod(key, period))}
    --- Subsection~5.1.
  \item Small-value direct lookup: \texttt{key < period} $\Rightarrow$
    \texttt{cycle(key) == cycle.values(key)} --- Subsection~5.2.
  \item Periodicity:
    $\operatorname{cycle}(\text{key})=
    \operatorname{cycle}(\text{key}+\text{period}\cdot m)$ for any number of
    loops --- Subsection~5.3.
  \item Multi-loop consistency: value at \texttt{key} is independent of which
    multiple of the period is added --- Subsection~5.4.
  \item Mod propagation: remainder modulo $d$ at any position equals remainder
    at the base position --- Subsection~5.5.
  \item Repeated-cycle invariance: repeating the base list preserves all
    lookups --- Subsection~5.6.
  \item Value positivity: non-negative base values guarantee non-negative cycle
    values --- Subsection~5.7.
  \item Rotation: rotating the base list shifts the cycle index by the same
    amount --- Subsection~5.8.
  \item MemCycle-level restatements: the key access and modulo properties are
    also verified directly for the memory-backed representation ---
    Subsection~5.9.
  \item Residue classification transfer: all-zero, none-zero, and some-zero
    divisor classifications on the base list hold at every cycle position ---
    Subsections~5.10--5.12.
\end{itemize}

Throughout this section, $L$ is the finite backing list, $n$ is the period,
and $Cycle$ is the unbounded periodic stream it defines --- formally,
\texttt{Cycle} means \texttt{RecCycle} (Subsection~3.1), proved in Section~4 to
equal \texttt{ModCycle} at every position:

\begin{equation*}
\begin{aligned}
\forall\,L\in\mathbb{L},\quad \forall\,v&\in\mathbb{N}_0,
  \quad \forall\,i\in\mathbb{N}_0 \\
L &:= [v_0,v_1,\dots,v_{n-1}]\in\mathbb{N}_0^n,\quad |L|>0 \\
\operatorname{Cycle} &:= [v_0,v_1,\dots,v_{n-1},v_0,v_1,\dots] \\
n &:= |L|
\end{aligned}
\end{equation*}

\subsection{Cycle Element Access}

The value of any element in a cycle equals the value of the underlying list at
the position modulo the cycle period.

\begin{equation*}
\operatorname{Cycle}_i=L[i\operatorname{mod}n]
\end{equation*}

\noindent\textbf{Proof.}

\begin{equation*}
\begin{aligned}
\operatorname{Cycle}_i &= \operatorname{RecCycle}_i
  && \text{[Subsection 3.1 Definition]} \\
&=\operatorname{ModCycle}_i && \text{[Section 4 Cycle Equivalence]} \\
&=L[i\operatorname{mod}n] && \text{[ModCycle Definition]}
\end{aligned}
\end{equation*}

\begin{equation*}
\therefore\quad \operatorname{Cycle}_i=L[i\operatorname{mod}n]
  \quad \blacksquare\ \text{[Q.E.D.]}
\end{equation*}

The Cycle Equivalence property was proved and verified in Section~4.

This property is verified in
\href{https://github.com/thiagomata/prime-numbers/blob/cycle-article-v1.0.0/src/main/scala/v1/chapter4/cycle/properties/CycleProperties.scala}{\texttt{CycleProperties::findValueInCycle}}.
The full Scala verification code is in Appendix~A.3.

\subsection{Small Value in Cycle}

For positions smaller than the cycle period, the cycle value equals the list
value at that position directly.

\begin{equation*}
i<n\implies\operatorname{Cycle}_i=L_i
\end{equation*}

\noindent\textbf{Proof.}

\begin{equation*}
\begin{aligned}
\operatorname{Cycle}_i &= \operatorname{RecCycle}_i
  && \text{[Subsection 3.1 Definition]} \\
i<n\implies\operatorname{RecCycle}_i &= L_i
  && \text{[RecCycle Definition]}
\end{aligned}
\end{equation*}

\begin{equation*}
\therefore\quad i<n\implies\operatorname{Cycle}_i=L_i
  \quad \blacksquare\ \text{[Q.E.D.]}
\end{equation*}

This step only needs the \texttt{Cycle := RecCycle} naming from
Subsection~3.1; the \texttt{ModCycle} side of the equivalence proved in
Section~4 is not required here.

This property is verified in
\href{https://github.com/thiagomata/prime-numbers/blob/cycle-article-v1.0.0/src/main/scala/v1/chapter4/cycle/properties/CycleProperties.scala}{\texttt{CycleProperties::smallValueInCycle}}.
The full Scala verification code is in Appendix~A.4.

\subsection{Value Match After Many Loops}

Cycle values remain invariant when adding any multiple of the cycle period to
the access key.

\begin{equation*}
\operatorname{Cycle}_{(i+n\cdot m)}=L[i\operatorname{mod}n]
\end{equation*}

\noindent\textbf{Proof.}

\begin{equation*}
\begin{aligned}
\operatorname{Cycle}_{(i+n\cdot m)}
  &=\operatorname{RecCycle}_{(i+n\cdot m)}
  && \text{[Subsection 3.1 Definition]} \\
  &=\operatorname{ModCycle}_{(i+n\cdot m)}
  && \text{[Section 4 Cycle Equivalence]} \\
  &=L[(i+n\cdot m)\operatorname{mod}n]
  && \text{[ModCycle Definition]} \\
  &=L[i\operatorname{mod}n]
  && \text{[Quotient Invariance Under Linear Shift by Multiplier]}
\end{aligned}
\end{equation*}

\begin{equation*}
\therefore\quad \operatorname{Cycle}_{(i+n\cdot m)}=L[i\operatorname{mod}n]
  \quad \blacksquare\ \text{[Q.E.D.]}
\end{equation*}

The lemma
\href{https://github.com/thiagomata/prime-numbers/blob/cycle-article-v1.0.0/articles/chapter2/modulo.md#65-quotient-invariance-under-linear-shift}{Quotient Invariance Under Linear Shift}
and its multiplier variant were proved and verified in
\emph{Division and Modulo from Recursive Normalization}~\cite{mata2026modulo}.

This property is verified in
\href{https://github.com/thiagomata/prime-numbers/blob/cycle-article-v1.0.0/src/main/scala/v1/chapter4/cycle/properties/CycleProperties.scala}{\texttt{CycleProperties::valueMatchAfterManyLoops}}.
The full Scala verification code is in Appendix~A.5.

\subsection{Two Multiples of Cycle Size}

Shifting the same key by two different multiples of the cycle period produces
the same cycle value either way.

\begin{equation*}
\operatorname{Cycle}_{(i+n\cdot m_1)}
  =\operatorname{Cycle}_{(i+n\cdot m_2)}
\end{equation*}

\noindent\textbf{Proof.} By Subsection~5.3, applied at both $m_1$ and $m_2$:

\begin{equation*}
\begin{aligned}
\operatorname{Cycle}_{(i+n\cdot m_1)}
  &=L[i\operatorname{mod}n]
  && \text{[Subsection 5.3, Value Match After Many Loops]} \\
\operatorname{Cycle}_{(i+n\cdot m_2)}
  &=L[i\operatorname{mod}n]
  && \text{[Subsection 5.3, Value Match After Many Loops]}
\end{aligned}
\end{equation*}

\begin{equation*}
\therefore\quad
\operatorname{Cycle}_{(i+n\cdot m_1)}
  =\operatorname{Cycle}_{(i+n\cdot m_2)}
  =L[i\operatorname{mod}n]
  \quad \blacksquare\ \text{[Q.E.D.]}
\end{equation*}

This property is verified in
\href{https://github.com/thiagomata/prime-numbers/blob/cycle-article-v1.0.0/src/main/scala/v1/chapter4/cycle/properties/CycleProperties.scala}{\texttt{CycleProperties::valueMatchAfterManyLoopsInBoth}}.
The full Scala verification code is in Appendix~A.6.

\subsection{Propagate Modulo from Value to Cycle}

The modulo operation applied to a cycle value can be equivalently applied to
the underlying list value at the modular index. Since the previous sections
prove that the cycle representations agree at every position, we can use the
modulo cycle definition directly: cycle lookup first reduces the position to
\texttt{i mod n}, and taking a remainder by any positive divisor \texttt{d}
preserves that same base-position reduction.

\begin{equation*}
\operatorname{Cycle}_i\operatorname{mod}d
  =L_{i\operatorname{mod}n}\operatorname{mod}d
\end{equation*}

\noindent\textbf{Proof.}

\begin{equation*}
\begin{aligned}
\operatorname{Cycle}_i
  &=\operatorname{RecCycle}_i && \text{[Subsection 3.1 Definition]} \\
  &=\operatorname{ModCycle}_i && \text{[Section 4 Cycle Equivalence]} \\
  &=L_{i\operatorname{mod}n} && \text{[ModCycle Definition]}
\end{aligned}
\end{equation*}

\begin{equation*}
\therefore\quad
\operatorname{Cycle}_i\operatorname{mod}d
  =L_{i\operatorname{mod}n}\operatorname{mod}d
  \quad \blacksquare\ \text{[Q.E.D.]}
\end{equation*}

This property is verified in
\href{https://github.com/thiagomata/prime-numbers/blob/cycle-article-v1.0.0/src/main/scala/v1/chapter4/cycle/properties/CycleProperties.scala}{\texttt{CycleProperties::propagateModFromValueToCycle}}.
The related idempotence restatement,
\texttt{cycle(position) == cycle(position mod period)}, is verified in
\href{https://github.com/thiagomata/prime-numbers/blob/cycle-article-v1.0.0/src/main/scala/v1/chapter4/cycle/properties/CycleProperties.scala}{\texttt{CycleProperties::assertCycleOfPosEqualsCycleOfModPos}}.
The full Scala verification code is in Appendix~A.7.

\subsection{Repeated-Cycle Invariance}

When a cycle's base list is repeated $t$ times to form a longer physical
period, the values read at every position remain identical. The repeated cycle
is structurally a concatenation of $t$ copies of the original list; the extra
length is invisible at any index because modular indexing composes correctly
across the nested periods. Repetition uses the same $t$-fold notation as
\href{https://github.com/thiagomata/prime-numbers/blob/cycle-article-v1.0.0/articles/chapter4/integral-cycle.md#61-integral-invariance-under-cycle-repetition}{\emph{Integral Cycle}, Subsection 6.1's $L^{(x)}$}
(list repeated $x$ times), defined the same way by unfolding one copy at a
time:

\begin{equation*}
\begin{aligned}
V^{(t)} &:={}
\begin{cases}
L_e & \text{if }t=0, \\
V\mathbin{\texttt{++}}V^{(t-1)} & \text{otherwise}.
\end{cases}
\end{aligned}
\end{equation*}

\begin{equation*}
\begin{aligned}
C &\text{ --- original MemCycle},\quad V=\operatorname{values}(C),
  \quad n=|V| \\
C^{(t)} &\text{ --- repeated cycle},\quad
  \operatorname{values}(C^{(t)})=V^{(t)}\quad\text{with }t>0 \\
\operatorname{period} &=t\cdot n
\end{aligned}
\end{equation*}

\noindent\textbf{Proof.}

\begin{equation*}
\begin{aligned}
C^{(t)}(\text{pos})
  &=V^{(t)}(\operatorname{mod}(\text{pos},\operatorname{period}))
  && \text{[MemCycle access via modular index]} \\
  &=V(\operatorname{mod}(\operatorname{mod}(\text{pos},t\cdot n),n))
  && \text{[Repeated list access pattern]} \\
  &=V(\operatorname{mod}(\text{pos},n))
  && \text{[ModOperations::modByPositiveMultipleThenBase]} \\
  &=C(\text{pos})
  && \text{[Original cycle access]} \\
  &\quad\blacksquare && \text{[Q.E.D.]}
\end{aligned}
\end{equation*}

The proof separates construction from lookup. Callers must build a valid
\texttt{MemCycle} from the repeated values; this lemma only says that once such
a cycle exists, the larger physical period does not change any lookup.

\begin{lstlisting}[style=scala]
def assertRepeatedValuesCycleMatches(
  cycle: MemCycle,
  repeatedCycle: MemCycle,
  times: BigInt,
  position: BigInt
): Boolean = {
  require(times > BigInt(0))
  require(position >= BigInt(0))
  require(cycle.period > BigInt(0))
  require(repeatedCycle.values == ListRepeatProperties.repeat(cycle.values, times))
  // ... inductive proof via mod composition ...
  repeatedCycle(position) == cycle(position)
}.holds
\end{lstlisting}

This property is verified in
\href{https://github.com/thiagomata/prime-numbers/blob/cycle-article-v1.0.0/src/main/scala/v1/chapter4/cycle/memory/properties/MemCycleProperties.scala}{\texttt{MemCycleProperties::assertRepeatedValuesCycleMatches}}.
The complete Scala derivation is included in Appendix~A.8.

\subsection{Cycle Value Positivity}

When every value in the base list is non-negative, every position in the cycle
returns a non-negative value. This guarantees that cycle lookups never produce
negative numbers, which is essential for integral and gap reasoning.

\begin{equation*}
\begin{aligned}
(\forall x\in L,\ x\geq 0)\land |L|>0
  \implies \operatorname{Cycle}_{\text{pos}}\geq 0
\end{aligned}
\end{equation*}

\noindent\textbf{Proof.} By Subsection~5.1,
$\operatorname{Cycle}_{\text{pos}}=L[\text{pos}\operatorname{mod}n]$, and
$\text{pos}\operatorname{mod}n$ is a valid index into $L$ (in $[0,n)$). So
$\operatorname{Cycle}_{\text{pos}}$ is one of the values of $L$---and every
value of $L$ is non-negative by hypothesis, so
$\operatorname{Cycle}_{\text{pos}}\geq 0$.

\begin{equation*}
\therefore\quad
(\forall x\in L,\ x\geq 0)\land |L|>0
  \implies \operatorname{Cycle}_{\text{pos}}\geq 0
  \quad \blacksquare\ \text{[Q.E.D.]}
\end{equation*}

This property is verified in
\href{https://github.com/thiagomata/prime-numbers/blob/cycle-article-v1.0.0/src/main/scala/v1/chapter4/cycle/properties/CycleProperties.scala}{\texttt{CycleProperties::cycleValuePositiveOrZero}},
which reduces to the list-level helper
\texttt{CycleUtils::checkPositiveOrZeroAtIndex}---indexing into a non-negative
list at a valid position yields a non-negative value. The full Scala
verification code is in Appendix~A.9.

The same argument gives the verified strict lower-bound form: if every base
value is greater than $x$, then every cycle value is greater than $x$.

\begin{equation*}
\begin{aligned}
(\forall y\in L,\ y>x)\land |L|>0
  &\implies \operatorname{Cycle}_{\text{pos}}>x.
\end{aligned}
\end{equation*}

\noindent\textbf{Proof.} The modular index selects a value of $L$, which is
greater than $x$ by hypothesis:

\begin{equation*}
\begin{aligned}
\operatorname{Cycle}_{\text{pos}}
  &=L[\text{pos}\operatorname{mod}n] && \text{[Subsection 5.1]} \\
L[\text{pos}\operatorname{mod}n]
  &>x && \text{[Base-list lower bound]} \\
\therefore\quad \operatorname{Cycle}_{\text{pos}}
  &>x.\quad \blacksquare\ \text{[Q.E.D.]}
\end{aligned}
\end{equation*}

{\raggedright
This strengthening is verified for \texttt{RecursiveCycle}, the article's
\texttt{Cycle} representation, in
\href{https://github.com/thiagomata/prime-numbers/blob/cycle-article-v1.0.0/src/main/scala/v1/chapter4/cycle/recursive/RecursiveCycle.scala#cycleValueBiggerThan}{\texttt{RecursiveCycle::cycleValueBiggerThan}}.
\par}

\subsection{Cycle Rotation}

Rotating a cycle's base list by $k$ positions and then accessing index $i$
gives the same value as accessing the original cycle at index $i+k$. This
connects cycle structure directly to the list rotation concept from chapter 3.
Rotation is defined by re-indexing the base list from offset $k$:

\begin{equation*}
\begin{aligned}
\operatorname{rotateAt}(L,k)
  &:=\big[\,L[(k+j)\operatorname{mod}n]\mid 0\leq j<n\,\big] \\
\operatorname{rotateAt}(\operatorname{Cycle},k)_i
  &:=\operatorname{rotateAt}(L,k)[i\operatorname{mod}n]
\end{aligned}
\end{equation*}

\begin{equation*}
\begin{aligned}
\operatorname{rotateAt}(\operatorname{Cycle},k)_i
  =\operatorname{Cycle}_{i+k}\quad\text{for }k\geq 0,\ i\geq 0
\end{aligned}
\end{equation*}

\noindent\textbf{Proof.}

\begin{equation*}
\begin{aligned}
\operatorname{rotateAt}(\operatorname{Cycle},k)_i
  &=\operatorname{rotateAt}(L,k)[i\operatorname{mod}n]
  && \text{[By Definition]} \\
  &=L[(k+(i\operatorname{mod}n))\operatorname{mod}n]
  && \substack{\text{[By Definition, since }i\operatorname{mod}n\\
                \text{ lies in }[0,n)\text{]}} \\
  &=L[(k+i)\operatorname{mod}n]
  && \substack{\text{[Modulo Idempotence + Distributivity}\\
                \text{over Addition]}} \\
  &=\operatorname{Cycle}_{i+k}
  && \text{[Subsection 5.1 Cycle Element Access]}
\end{aligned}
\end{equation*}

\begin{equation*}
\therefore\quad
\operatorname{rotateAt}(\operatorname{Cycle},k)_i
  =\operatorname{Cycle}_{i+k}\quad\text{for }k\geq 0,\ i\geq 0
  \quad \blacksquare\ \text{[Q.E.D.]}
\end{equation*}

The third step composes
\href{https://github.com/thiagomata/prime-numbers/blob/cycle-article-v1.0.0/articles/chapter2/modulo.md#68-modulo-idempotence}{Modulo Idempotence}
and
\href{https://github.com/thiagomata/prime-numbers/blob/cycle-article-v1.0.0/articles/chapter2/modulo.md#69-distributivity-over-addition}{Distributivity over Addition},
both proved and verified in
\emph{Division and Modulo from Recursive Normalization}~\cite{mata2026modulo}:
$\operatorname{mod}(k+\operatorname{mod}(i,n),n)
=\operatorname{mod}(k+i,n)$ follows because both sides reduce to
$\operatorname{mod}(\operatorname{mod}(k,n)+
\operatorname{mod}(i,n),n)$.

This property is verified in
\href{https://github.com/thiagomata/prime-numbers/blob/cycle-article-v1.0.0/src/main/scala/v1/chapter4/cycle/properties/CycleProperties.scala}{\texttt{CycleProperties::rotateAtValue}}.
The full Scala verification code is in Appendix~A.10.

\subsection{MemCycle-Level Restatement}

Sections~5.1--5.5 state element access, small-value lookup, periodic
invariance, multi-loop consistency, and mod propagation for \texttt{ModCycle}.
Since $\operatorname{MemCycle}(L)_i=\operatorname{ModCycle}(L)_i$ at every
position (Subsection~3.3, assembled at the close of Section~4), all five
results carry over to \texttt{MemCycle} by direct substitution---mathematically,
this section adds no content beyond Subsections~5.1--5.5. \texttt{MemCycle} is
nonetheless the representation actually used elsewhere in the codebase (it is
the one that carries residue-classification metadata), and Stainless has no way
to transfer a lemma proved for one concrete type to a different wrapper type
that merely happens to agree with it pointwise: each result below is re-proved
against \texttt{MemCycle}, and its proof body mirrors the corresponding
\texttt{ModCycle} lemma line for line rather than deriving anything new:

\begin{equation*}
\begin{aligned}
\operatorname{Cycle}_{\text{key}}
  &=L[\text{key}\operatorname{mod}n] && \text{[Element Access]} \\
\text{key}<n
  &\implies \operatorname{Cycle}_{\text{key}}=L[\text{key}]
  && \text{[Small Value Lookup]} \\
\operatorname{Cycle}_{\text{key}}
  &=\operatorname{Cycle}_{\text{key}+n\cdot m}
  && \text{[Periodic Invariance]} \\
\operatorname{Cycle}_{\text{key}+n\cdot m_1}
  &=\operatorname{Cycle}_{\text{key}+n\cdot m_2}
  && \text{[Multi-Loop Consistency]} \\
\operatorname{Cycle}_{\text{key}}\operatorname{mod}d
  &=L[\text{key}\operatorname{mod}n]\operatorname{mod}d
  && \text{[Mod Propagation]}
\end{aligned}
\end{equation*}

{\raggedright
These are verified in
\href{https://github.com/thiagomata/prime-numbers/blob/cycle-article-v1.0.0/src/main/scala/v1/chapter4/cycle/memory/properties/MemCycleProperties.scala}{\texttt{MemCycleProperties::findValueInCycle}},
\href{https://github.com/thiagomata/prime-numbers/blob/cycle-article-v1.0.0/src/main/scala/v1/chapter4/cycle/memory/properties/MemCycleProperties.scala}{\texttt{MemCycleProperties::smallValueInCycle}},
\href{https://github.com/thiagomata/prime-numbers/blob/cycle-article-v1.0.0/src/main/scala/v1/chapter4/cycle/memory/properties/MemCycleProperties.scala}{\texttt{MemCycleProperties::valueMatchAfterManyLoops}},
\href{https://github.com/thiagomata/prime-numbers/blob/cycle-article-v1.0.0/src/main/scala/v1/chapter4/cycle/memory/properties/MemCycleProperties.scala}{\texttt{MemCycleProperties::valueMatchAfterManyLoopsInBoth}},
and
\href{https://github.com/thiagomata/prime-numbers/blob/cycle-article-v1.0.0/src/main/scala/v1/chapter4/cycle/memory/properties/MemCycleProperties.scala}{\texttt{MemCycleProperties::propagateModFromValueToCycle}}.

The mod-idempotence identity from Subsection~5.5's proof (that
$\operatorname{Cycle}_i$ equals
$\operatorname{Cycle}_{(i\operatorname{mod}n)\operatorname{mod}n}$) has its
own \texttt{MemCycle} restatement in
\href{https://github.com/thiagomata/prime-numbers/blob/cycle-article-v1.0.0/src/main/scala/v1/chapter4/cycle/memory/properties/MemCycleProperties.scala}{\texttt{MemCycleProperties::assertCycleOfPosEqualsCycleOfModPos}}.
\par}

{\raggedright
\texttt{MemCycle} also classifies each divisor by how its residues behave
across the base list. \texttt{MemCycle.checkMod(d)} (Subsection~3.3) counts how
many values of $L$ are $\equiv 0\operatorname{mod}d$, via
\texttt{countModZero}---one pass over the list, adding 1 wherever the value is
$\equiv 0\operatorname{mod}d$:
\par}

\begin{equation*}
\begin{aligned}
\operatorname{countModZero}(L,d)&:={}
\begin{cases}
0 & \text{if }L=L_e, \\
1+\operatorname{countModZero}(\operatorname{tail}(L),d)
  & \text{if }\operatorname{head}(L)\operatorname{mod}d=0, \\
\operatorname{countModZero}(\operatorname{tail}(L),d)
  & \text{otherwise}.
\end{cases}
\end{aligned}
\end{equation*}

\texttt{checkMod(d)} then buckets $d$ as all-zero, none-zero, or some-zero
based on that count---each bucket defined and proved below to transfer from the
base list to every position of the infinite \texttt{Cycle}, using the same
move as Subsection~5.7: a cycle value is always \emph{some} value of $L$
(Subsection~5.1), so any property that holds of every value of $L$ holds of
every cycle position too. \texttt{CycleCheckMod.scala}'s ten lemmas separately
verify that the classification itself is mutually exclusive, exhaustive, and
persists correctly across independent \texttt{checkMod} calls---bookkeeping
properties of the classification lists, not further claims about cycle values,
so they are not restated here.

\subsection{All-Zero Residue Transfers to the Cycle}

\texttt{checkMod(d)} classifies \texttt{d} as all-zero when every value of
\texttt{L} is $\equiv 0\operatorname{mod}d$---i.e.
\texttt{countModZero(L, d)} counts all $n$ of them:

\begin{equation*}
\begin{aligned}
\operatorname{allZero}(d)
  &:\Leftrightarrow \operatorname{countModZero}(L,d)=n
  && \text{[MemCycle.allModValuesAreZero]}
\end{aligned}
\end{equation*}

\begin{equation*}
\begin{aligned}
\operatorname{allZero}(d)
  \implies \forall k,\ \operatorname{Cycle}_k\operatorname{mod}d=0
\end{aligned}
\end{equation*}

\noindent\textbf{Proof.}

\begin{equation*}
\begin{aligned}
\operatorname{allZero}(d)
  &\implies \forall x\in L,\ x\operatorname{mod}d=0
  && \text{[Definition of countModZero]} \\
\operatorname{Cycle}_k
  &=L[k\operatorname{mod}n]
  && \text{[Subsection 5.1 Cycle Element Access]} \\
\implies \operatorname{Cycle}_k\operatorname{mod}d
  &=0
  && \text{[$L[k\operatorname{mod}n]$ is a value of $L$]}
\end{aligned}
\end{equation*}

\begin{equation*}
\therefore\quad
\operatorname{allZero}(d)
  \implies \forall k,\ \operatorname{Cycle}_k\operatorname{mod}d=0
  \quad \blacksquare\ \text{[Q.E.D.]}
\end{equation*}

\texttt{allZero} is \texttt{MemCycle.allModValuesAreZero}, defined and set by
\texttt{checkMod} in
\href{https://github.com/thiagomata/prime-numbers/blob/cycle-article-v1.0.0/src/main/scala/v1/chapter4/cycle/memory/MemCycle.scala}{\texttt{MemCycle}}.
Stainless verifies the base-list count directly; the transfer to an arbitrary
cycle position \texttt{k} shown above is this article's corollary of that count
together with Subsection~5.1, not a separate Stainless lemma.

\subsection{None-Zero Residue Transfers to the Cycle}

The symmetric case: \texttt{checkMod(d)} classifies \texttt{d} as none-zero
when no value of \texttt{L} is $\equiv 0\operatorname{mod}d$---
\texttt{countModZero(L, d)} counts none of them:

\begin{equation*}
\begin{aligned}
\operatorname{noneZero}(d)
  &:\Leftrightarrow \operatorname{countModZero}(L,d)=0
  && \text{[MemCycle.noModValuesAreZero]}
\end{aligned}
\end{equation*}

\begin{equation*}
\begin{aligned}
\operatorname{noneZero}(d)
  \implies \forall k,\ \operatorname{Cycle}_k\operatorname{mod}d\neq 0
\end{aligned}
\end{equation*}

\noindent\textbf{Proof.}

\begin{equation*}
\begin{aligned}
\operatorname{noneZero}(d)
  &\implies \forall x\in L,\ x\operatorname{mod}d\neq 0
  && \text{[Definition of countModZero]} \\
\operatorname{Cycle}_k
  &=L[k\operatorname{mod}n]
  && \text{[Subsection 5.1 Cycle Element Access]} \\
\implies \operatorname{Cycle}_k\operatorname{mod}d
  &\neq 0
  && \text{[$L[k\operatorname{mod}n]$ is a value of $L$]}
\end{aligned}
\end{equation*}

\begin{equation*}
\therefore\quad
\operatorname{noneZero}(d)
  \implies \forall k,\ \operatorname{Cycle}_k\operatorname{mod}d\neq 0
  \quad \blacksquare\ \text{[Q.E.D.]}
\end{equation*}

\texttt{noneZero} is \texttt{MemCycle.noModValuesAreZero}, defined and set by
\texttt{checkMod} in
\href{https://github.com/thiagomata/prime-numbers/blob/cycle-article-v1.0.0/src/main/scala/v1/chapter4/cycle/memory/MemCycle.scala}{\texttt{MemCycle}}.
As in Subsection~5.10, Stainless verifies the base-list count; the transfer to
arbitrary \texttt{k} is this article's corollary of Subsection~5.1.

\subsection{Some-Zero Residue Transfers to the Cycle}

\texttt{checkMod(d)} classifies \texttt{d} as some-zero when some, but not all,
values of \texttt{L} are $\equiv 0\operatorname{mod}d$---
\texttt{countModZero(L, d)} counts strictly between \texttt{0} and $n$ of them:

\begin{equation*}
\begin{aligned}
\operatorname{someZero}(d)
  &:\Leftrightarrow 0<\operatorname{countModZero}(L,d)<n
  && \text{[MemCycle.someModValuesAreZero]}
\end{aligned}
\end{equation*}

Unlike Subsections~5.10--5.11, \texttt{someZero} is not a claim about
\emph{every} position---it needs witnesses on both sides, so the transfer
argument has to name two concrete cycle positions rather than substitute into
a universal.

\begin{equation*}
\begin{aligned}
\operatorname{someZero}(d)\implies \exists\,k_0,k_1,\quad
&\operatorname{Cycle}_{k_0}\operatorname{mod}d=0\ \land{} \\
&\operatorname{Cycle}_{k_1}\operatorname{mod}d\neq 0
\end{aligned}
\end{equation*}

\noindent\textbf{Proof.} By definition, \texttt{someZero(d)} means
$0<\operatorname{countModZero}(L,d)<n$: at least one value of \texttt{L} is
$\equiv 0\operatorname{mod}d$, and at least one is not. Let
$j_0,j_1\in[0,n)$ be indices of one such value on each side, so
$L_{j_0}\operatorname{mod}d=0$ and
$L_{j_1}\operatorname{mod}d\neq 0$.

\begin{equation*}
\begin{aligned}
j_0,j_1<n
  &\implies \operatorname{Cycle}_{j_0}=L_{j_0}\ \land\
    \operatorname{Cycle}_{j_1}=L_{j_1}
  && \text{[Subsection 5.2 Small Value in Cycle]} \\
  &\implies \operatorname{Cycle}_{j_0}\operatorname{mod}d=0\ \land\
    \operatorname{Cycle}_{j_1}\operatorname{mod}d\neq 0
  && \text{[Substitution]}
\end{aligned}
\end{equation*}

\begin{equation*}
\begin{gathered}
\therefore\quad
\operatorname{someZero}(d)\implies \exists\,k_0,k_1,\quad
\operatorname{Cycle}_{k_0}\operatorname{mod}d=0\ \land\
\operatorname{Cycle}_{k_1}\operatorname{mod}d\neq 0 \\
(k_0:=j_0,\ k_1:=j_1)\quad \blacksquare\ \text{[Q.E.D.]}
\end{gathered}
\end{equation*}

\texttt{someZero} is \texttt{MemCycle.someModValuesAreZero}, defined and set
by \texttt{checkMod} in
\href{https://github.com/thiagomata/prime-numbers/blob/cycle-article-v1.0.0/src/main/scala/v1/chapter4/cycle/memory/MemCycle.scala}{\texttt{MemCycle}}.
As above, Stainless verifies the base-list count; the witness positions
$k_0,k_1$ are this article's corollary of Subsection~5.2, not a separate
Stainless lemma.
